\section{相关工作}
\label{sec:related_work}

\subsection{自主实验室中的仿真训练}
机器人自动化的集成显著加速了材料科学和化学领域的实验吞吐量。诸如移动机器人化学家\cite{burger2020mobile}和自动驾驶薄膜实验室\cite{macleod2020self}等开创性工作证明了自动化执行的潜力。然而，这些系统主要依赖于硬编码的路径点。物理仿真工具支持在实际部署前训练和测试接触密集型机器人技能。复杂的生物分析需要高保真的接触物理特性，而不仅仅是运动学边界框近似。为了解决这个问题，我们的框架采用了GPU加速的NVIDIA Isaac Sim\cite{makoviychuk2021isaac}，为模拟接触密集型操作提供了数学上严谨的基于张量的先验知识。

\subsection{安全导航与防流体溢出}
在处理挥发性试剂或生物培养物的安全关键型实验室环境中，确保绝对无碰撞和防溢出的移动底盘导航至关重要。诸如$RRT^{*}$或$A^{*}$等标准的基于采样的规划器擅长寻找无障碍路径，但通常忽略任务空间的方向约束，导致它们不适合运输敞口的液体容器。为了保证绝对的流体稳定性，我们将工作站之间的宏观转移公式化为一个有约束的经典优化问题。通过对末端执行器施加严格的滚转-俯仰边界并最小化轨迹加加速度积分，我们的经典规划模块在物理部署之前就系统地消除了液体溢出的风险。

\subsection{面向灵巧抓取的大规模并行强化学习}
虽然经典规划保证了宏观安全性，但为灵巧抓取和实验器材操作明确建立高度非线性的微观接触动力学模型在计算上是不可行的。深度强化学习（DRL）在掌握接触密集型任务方面表现出巨大潜力\cite{johannink2019residual}。然而，在现实世界中训练稳健的DRL策略受到样本效率低下和硬件磨损的限制。

为了克服这一问题，我们的框架利用Isaac Sim高度并行化的张量架构，在数千个并发虚拟环境中训练近端策略优化（PPO）抓取智能体。通过设计姿态感知奖励函数——在靠近和抬起阶段严厉惩罚方向偏差——策略学会了稳固地抓取试剂瓶而不会使其倾斜。结合广泛的域随机化，这种大规模并行RL管线稳健地弥合了模拟与现实之间的差距，成功部署到了真实的UR3机械臂上。
